The Tinker-\/\-H\-P suite is made of 5 main programs that have been designed to run in a High Performance Computing (H\-P\-C) context, on G\-P\-Us and C\-P\-Us, leveraging M\-P\-I to run with multiple processes. All of these run based on various {\bfseries input files} and a {\bfseries command line}.

Some input files are common to all the binaries\-:


\begin{DoxyItemize}
\item An {\bfseries xyz} file in the Tinker format representing a {\bfseries geometry} (through cartesian coordinates), {\bfseries atom types} of the atoms of the systems as well as their {\bfseries connectivity}.

{\itshape Below is an example of an xyz file for a two water molecule system, oxygens are of atom type 101, hydrogens of type 88.}
\end{DoxyItemize}


\begin{DoxyCode}
1  O      8.679662    7.087692   -0.696862    101     2     3
2  H      7.809455    6.755792   -0.382259     88     1
3  H      8.722232    6.814243   -1.617561     88     1
4  O     -0.117313    8.244447    6.837616    101     5     6
5  H      0.216892    7.895445    6.050027     88     4
6  H      0.444268    7.826013    7.530196     88     4
\end{DoxyCode}



\begin{DoxyItemize}
\item A {\bfseries keyfile} containing various keywords describing the simulation such as the size of the box, the cutoffs on non-\/bonded simulations, the integrator for Molecular Dynamics, and also the {\bfseries parameter file}.

{\itshape Below is an example of a keyfile for a system using the charmm22 parameter file, with a cubic box size of edge 18.\-643 Angstroms and the verlet integrator.}
\end{DoxyItemize}


\begin{DoxyCode}
parameters        charmm22
verbose
a-axis            18.643
integrator \hyperlink{verlet_8f90_adc85925ca8af0efc52d80fe31a74bfaa}{verlet}
\end{DoxyCode}



\begin{DoxyItemize}
\item A {\bfseries parameter file} listing the parameters, for each atom type, of the force field used for the simulation.
\end{DoxyItemize}\hypertarget{md_binaries_doxygen_running_the_programs}{}\section{Running the Programs}\label{md_binaries_doxygen_running_the_programs}
For all programs of the Tinker-\/\-H\-P suite, the geometry needs to be included in the command line. To further indicate to the programs that the simulation needs to take into account a specific keyfile, the keyfile needs to have the same prefix as the geometry or can be explicitly given in the command line with the prefix {\ttfamily -\/k keyfile}.

{\itshape Example of a command line for the analyze binary involving the {\ttfamily watersmall.\-xyz} geometry and the {\ttfamily watersmall.\-key} keyfile, on a single M\-P\-I process\-:}


\begin{DoxyCode}
mpirun -np 1 /path/to/\hyperlink{analyze_8f90_ad0571e0cf34d424af56a4981e316f8eb}{analyze} watersmall.xyz -k watersmall.key e
\end{DoxyCode}


The 5 programs are listed below with a description of the command line necessary to run them. In the case of an incomplete command line or a missing input file, an error message describing the missing point will be printed.\hypertarget{md_binaries_doxygen_program_analyze}{}\section{analyze}\label{md_binaries_doxygen_program_analyze}
Run a single point energy given a trajectory file or a single frame.
\begin{DoxyItemize}
\item {\bfseries Command line\-:} {\ttfamily E} as in \char`\"{}\-Energy decomposition analysis\char`\"{} (see example above)
\end{DoxyItemize}\hypertarget{md_binaries_doxygen_program_dynamic}{}\section{dynamic}\label{md_binaries_doxygen_program_dynamic}
Runs a molecular dynamics trajectory.
\begin{DoxyItemize}
\item {\bfseries Command line\-:} A list of arguments has to be passed directly through the command line, respectively\-:
\begin{DoxyItemize}
\item Number of steps of the desired trajectory
\item Duration of a timestep (in femtoseconds)
\item Frequency of output of frames (in picoseconds)
\item Integer corresponding to the statistical ensemble to be sampled\-:
\begin{DoxyItemize}
\item {\ttfamily 1}\-: N\-V\-E ensemble
\item {\ttfamily 2}\-: N\-V\-T ensemble
\item {\ttfamily 4}\-: N\-P\-T ensemble
\end{DoxyItemize}
\item Temperature (in Kelvin) for N\-V\-T and N\-P\-T ensembles
\item Pressure (in Atmosphere) for N\-P\-T ensemble
\end{DoxyItemize}
\end{DoxyItemize}

{\itshape Example of a command line for the dynamic binary corresponding to 10000 time steps of 1fs in the N\-P\-T ensemble at 300\-K and 1 Atm, with frames printed out every picosecond\-:}


\begin{DoxyCode}
mpirun -np 1 /path/to/\hyperlink{dynamic_8f90_a242c646ad028ae81e8a4bed2193eae5d}{dynamic} watersmall.xyz 10000 1 1 4 300 1
\end{DoxyCode}


\begin{quotation}
{\bfseries Note\-:} More information about the features available with the {\bfseries dynamic} program, such as integrators, thermostats and barostats, can be found \href{md_dynamic_doxygen.html}{\tt here}.

\end{quotation}
\hypertarget{md_binaries_doxygen_program_minimize}{}\section{minimize}\label{md_binaries_doxygen_program_minimize}
Runs a minimization of the potential energy of a system, starting from a geometry.
\begin{DoxyItemize}
\item {\bfseries Command line\-:} Convergence criteria of the minimization algorithm
\end{DoxyItemize}

{\itshape Example of a command line for the minimize binary corresponding to a convergence threshold of 1 for R\-M\-S gradient\-:}


\begin{DoxyCode}
mpirun -np 1 /path/to/\hyperlink{minimize_8f90_a0ea869a84e4dbfc9951ec265abef5b8d}{minimize} watersmall.xyz 1
\end{DoxyCode}


\begin{quotation}
{\bfseries Note\-:} The only optimizer is the limited memory B\-F\-G\-S algorithm so no information about the optimizer has to be written in the keyfile. Starting from a $\ast$.xyz file, the program will output an addition $\ast$.xyz\-\_\-2 file corresponding to the minimized geometry.

\end{quotation}
\hypertarget{md_binaries_doxygen_program_pimd}{}\section{P\-I\-M\-D}\label{md_binaries_doxygen_program_pimd}
Runs a Path Integral Molecular Dynamics trajectory to include nuclear quantum effects in the statistics.
\begin{DoxyItemize}
\item {\bfseries Command line\-:} The exact same inputs as the one of the {\bfseries dynamic} programs are necessary in the command line as described above. Several features specific to Path Integral such as the number of beads, as well as other implemented techniques to include nuclear quantum effects (the adaptive Quantum Thermal Bath) can be found in {\ttfamily Quantum-\/\-H\-P.\-md}.
\end{DoxyItemize}\hypertarget{md_binaries_doxygen_program_bar}{}\section{bar}\label{md_binaries_doxygen_program_bar}
Compute free energy differences through the Bennett Acceptance Ratio estimator. In Tinker-\/\-H\-P, free energy differences are obtained in a post-\/processing step, in two stages. Given two keyfiles and two trajectories associated to two hamiltonians, first a {\ttfamily .bar} file is produced with, for each trajectory, the potential energies associated to the two hamiltonians. Then, given the {\ttfamily .bar} file, one can run the self-\/consistent procedure to get the final free energy difference through the bar estimator.


\begin{DoxyItemize}
\item {\bfseries Command line\-:}
\begin{DoxyItemize}
\item {\ttfamily 1} to generate the {\ttfamily .bar} file, {\ttfamily 2} to post-\/process it to get the bar estimation.
\item For generation of the bar file, one has to indicate the two geometries with their keyfiles and the two temperatures.
\end{DoxyItemize}
\end{DoxyItemize}

{\itshape Example of a command line for the bar binary to generate the {\ttfamily .bar} file corresponding to systema at temperature tempa and systemb at temperature tempb\-:}


\begin{DoxyCode}
mpirun -np 1 /path/to/\hyperlink{bar_8f90_a75dfe3f0be9db5c03a0c04398460d0b5}{bar} 1 systema.xyz -k systema.key 300 systemb.xyz -k systemb.key 300
\end{DoxyCode}



\begin{DoxyItemize}
\item For estimation of the free energy difference, one has to indicate the bar file.
\end{DoxyItemize}

{\itshape Example of a command line for the bar binary to get the free energy estimation based on the {\ttfamily system.\-bar} file\-:}


\begin{DoxyCode}
mpirun -np 1 /path/to/\hyperlink{bar_8f90_a75dfe3f0be9db5c03a0c04398460d0b5}{bar} 2 system.bar
\end{DoxyCode}
\hypertarget{md_binaries_doxygen_program_testgrad}{}\section{testgrad}\label{md_binaries_doxygen_program_testgrad}
Compute gradients associated with a geometry. One can compute analytical gradients of the potential energy terms or numerical ones, on all the cartesian degrees of freedom.
\begin{DoxyItemize}
\item {\bfseries Command line\-:}
\begin{DoxyItemize}
\item {\ttfamily Y} to compute analytical gradients, {\ttfamily N} if it is not the case
\item {\ttfamily Y} to compute numerical gradients (through finite differences), {\ttfamily N} if not
\item When numerical gradients are requested, size of the step displacement in Angstroms
\item {\ttfamily Y} if breakdown by component is requested, {\ttfamily N} if not 


\end{DoxyItemize}
\end{DoxyItemize}\hypertarget{md_binaries_doxygen_general_information}{}\section{General Information}\label{md_binaries_doxygen_general_information}

\begin{DoxyItemize}
\item Tinker-\/\-H\-P is compatible with numerous bonded and non-\/bonded potential energy terms. For each of them, energies and analyticals gradients can be computed by running the programs described above. Each of them is associated with a number of features that can be modified by some keywords to be added to the keyfile. A list of these potential energy terms with associated keywords can be found \href{md_potential_doxygen.html}{\tt here}.
\item Informations about I/\-O can be found \href{md_io_doxygen.html}{\tt here}
\item Info about free energy simulations can be found \href{md_free_energy.html}{\tt here}, with specifics about alchemical ones through Lambda-\/\-A\-B\-F and related methodologies \href{md_lambda-abf_doxygen.html}{\tt here} 
\end{DoxyItemize}